% Hand-written sentences that depend on results; every number is a macro from v15_numbers.tex (Research/code/v15_tables.py).
\newcommand{\nAbstractModel}{Model-based analyses trade this robustness for efficiency: a log-linear working model fitted to the same capped records is more efficient at a thousand records but biased, and imputation from complete source records loses coverage under pre-specified population shifts.}
\newcommand{\nPointerSentence}{at $b\in\{3,4\}$ fixed forms lose exactly $4$ and $8$, two rounds lose $1.80$ and $3.32$--$3.33$ against the upper bounds $2$ and $4$ of part (a), and $b-1$ rounds lose $1.80$ at both caps (Appendix~\ref{app:pointerexact}).}
\newcommand{\nDiscussionSentence}{optimized forms needed between $1.0$ and \nMaxAll{} times the records of branching, where balanced random forms needed an order of magnitude more; one adaptive round recovered at least \nTwoMinPct\% of the log-gap in every case solved to optimality; label-dependent form probabilities did not change the picture; and designs fitted on reference records lost a few percent while remaining exactly unbiased.}
\newcommand{\nOrderSentence}{in the position model of Section~\ref{sec:experiments}, a 10\% chance that the last answer flips biases branching and form scores by a median \nPosBranch{} and \nPosFit{} branching standard errors at $n=1000$.}
\newcommand{\nModelParagraph}{The constants above are for design-unbiased estimation; a model can buy efficiency, at a price. \emph{Complete source records.} Imputing the unasked items from complete records of another population uses information that the capped audit does not collect. When the conditional law of the unasked items transports (random-half and sex splits), regression and multiple imputation attain RMSE \nImpRmseLo--\nImpRmseHi{} times that of branching with coverage \nImpGoodCovLo--\nImpGoodCovHi. In \nImpBadN{} of the eight settings with a pre-specified shift (education at both caps, poor health at $b=5$), regression imputation is biased by up to \nImpBiasHi{} branching standard deviations and covers \nImpCovLo--\nImpCovHi, multiple imputation covers \nMiCovLo--\nMiCovHi, and the design-based scores cover \nDesCovLo--\nDesCovHi{} in every setting (Table~\ref{tab:shift}). \emph{A working model on the capped records.} A log-linear model with two-way interactions among the label and the rule's items, fitted by EM to the records of the fitted forms, needs no reference \citep{schafer1997}. At $n=1000$ its RMSE is \nEmRmseLo--\nEmRmseHi{} times that of the design-based form score, on a par with branching, with coverage \nEmCovLo--\nEmCovHi. Its bias does not shrink with $n$: at $n=\nEmBigN$ its coverage is \nEmBigCovLo--\nEmBigCovHi{} and its RMSE exceeds that of the form score in \nEmBigWorse{} of \nEmBigSettings{} settings, while branching and the form score keep coverage \nEmBigFixCovLo--\nEmBigFixCovHi{} (Table~\ref{tab:em}). On balanced random forms the same model has RMSE \nEmUnifLo--\nEmUnifHi{} times that of the optimized forms\textquotesingle{} design-based score at $n=1000$. The working model is thus a strong competitor at survey sizes of a thousand, and its error is governed by a bias that the capped records cannot bound; design-based scores are the tool when a claim must not rest on a model, and that is the regime in which the constants of this paper price pre-commitment.}
\newcommand{\nEmAppendixText}{At $n=1000$ the working model on the fitted-form records has the smallest RMSE of the four estimators in most settings, \nEmBrLo--\nEmBrHi{} times that of branching: the model shares strength across paths that the design-based scores treat separately. Its bias is between \nEmBiasLo{} and \nEmBiasHi{} branching standard deviations and its coverage \nEmCovLo--\nEmCovHi{} (below 0.925 in \nEmBelow{} of \nEmSettings{} settings). The bias is a property of the model, not of the sample size: the accuracy of a depth-$b$ tree depends on interactions of up to $b$ items and the label, which a two-way model does not represent. At $n=\nEmBigN$ the same bias is up to \nEmBigBiasHi{} of the estimator\textquotesingle s own standard deviations, coverage falls to \nEmBigCovLo--\nEmBigCovHi{} (below 0.90 in \nEmBigBelow{} of \nEmBigSettings{} settings), and the RMSE is \nEmBigRmseLo--\nEmBigRmseHi{} times that of branching and above that of the design-based form score in \nEmBigWorse{} settings; the design-based scores keep coverage \nEmBigFixCovLo--\nEmBigFixCovHi. A richer model would reduce the bias and raise the variance; with all interactions it is the design-based estimator. On balanced random forms the model-based estimate is unbiased to a similar degree but \nEmUnifLo--\nEmUnifHi{} times less accurate at $n=1000$ than the design-based score on optimized forms, so the choice of forms matters for model-based analysis as well.}
